\subsection{Against the default approach}\label{s:missed}

A classifier, like the signature check, works only if its fit covers the declared work.
\Cref{t:t6} sets $s_R$ beside the three classifier fits. It lists what each detector
learns from, its detection of the unthrottled co-location run at $\tau_\alpha$, and its
false-positive rate on test decode.

\begin{table*}[!tp]
  \centering
  \small
  % generated by scripts/table_t6.py from (repo_commit of each source):
% operating_point.json 1d2e0a2, h2.json 41b1fed, r1_scores.json 53f9463,
% windows/manifest.json 018c3a2, replication/t3.json (amodo 73ffebe)
\begin{tabular}{lccccrrl}
\toprule
 & \multicolumn{4}{c}{learns from} & & & \\
\cmidrule(lr){2-5}
detector & declared & honest training & evasions & non-AI & $d_{\mathrm{unthrottled}}(\tau_\alpha)$ & FPR decode & alarm means \\
\midrule
replica & pub. & pub. & -- & pub. & 0.6\% & 0.0\% & training \\
no-evasion fit & \checkmark & \checkmark & -- & \checkmark & 100.0\% & 5.3\% & training \\
$s_C^{\mathrm{dev}}$ & \checkmark & \checkmark & \checkmark & \checkmark & 100.0\% & 1.5\% & training \\
$s_R$ & \checkmark & -- & -- & -- & 100.0\% & 0.6\% & undeclared \\
\bottomrule
\end{tabular}

  \caption{The data each detector learns from, its detection of the unthrottled run at
    $\tau_\alpha$, and its false-positive rate on test decode. ``pub.'': traces recorded
    on the same GPU and probe to the published design, not among the \NCaptures{}
    runs of \cref{s:corpus}. An alarm of $s_R$ means ``undeclared''; an alarm of a
    classifier means ``training''. The rows for the replica and the no-evasion fit are
    exploratory.}
  \label{t:t6}
\end{table*}

At its $\tau_\alpha$, the classifier, $s_C^{\mathrm{dev}}$, does nearly as well as $s_R$.
It detects $\Xprime\%$ of co-location, and every window of every other group except
jitter, where its weakest run is detected at \ClfJitterWorst\%. It stands for the
default approach because it learns from examples of training. But it is fitted on the
monitored GPU, declared work included, as the signature check is, so it needs no
transfer between chips. It is therefore the most favourable case of the default
approach. Even without evasion examples, as in the no-evasion fit, a classifier
needs examples of honest training, whereas $s_R$ needs only recordings of the declared
work.

The threshold of $s_C^{\mathrm{dev}}$ is also fragile. The classifier scores
\FracCalZero\% of calibration declared windows exactly zero, so
$\tau_\alpha = \TauClf$. It then flags any window that stage one calls AI and that even
one of its \ForestTrees{} stage-two trees votes as training. Small changes in declared
work could therefore shift its flags.

A classifier gives no guarantee on a declared workload that its fit has not seen. The
replica, fitted to traces recorded to the published design, detects only
\XzeroPrime\% of the unthrottled run at its $\tau_\alpha$, and flags \RepPrefill\% of
development prefill. $s_C^{\mathrm{dev}}$, which has development prefill in its fit,
flags \ClfDevPrefill\% of it. \citet{gargiulo2026} publish two forests fitted to their
own corpus: a genuine forest, fitted to their training, inference and non-AI traces,
and a hardened forest, which adds their adversarial traces. The two disagree on our
windows (\cref{t:t5}). At their default threshold of \DefaultDecision, the hardened
forest catches all three co-location runs, but it flags \GkHardPrefill\% of declared
prefill windows as training. The genuine forest flags only \GkGenPrefill\% of prefill,
but it misses the unthrottled run ($\GkGenMiss\%$). Nothing in either fit tells the
verifier which outcome to expect. The replica's traces differ from ours in workload
design and session, so its results do not rank the signature check against the default
approach.
